% ============================== CI-2 ==============================
\reading{CI-2}{The Financial Stability Implications of Artificial Intelligence}{STRIKE FIRST}{6}

\begin{cikeybox}[title={The three objectives, in the spine's own words}]
\begin{enumerate}[label=\textbf{\alph*.}]
\item Explain key developments in AI since 2017, differentiating between supply-side drivers
and demand-side drivers.
\item Identify and analyze existing and emerging AI use cases in financial services,
including industry use cases as well as regulatory and supervisory use cases.
\item Assess the potential implications of AI for financial stability, with a focus on the
following vulnerabilities: third-party dependencies and service provider concentration;
market correlations; cyber risks; model risk, data, and governance; other vulnerabilities,
such as fraud, disinformation, and misalignment.
\end{enumerate}
\greyline{Source: ``The Financial Stability Implications of Artificial Intelligence'', FSB
(November 2024). The class notes label this Reading 101; it is CI-2 on the 2026 spine.
Objective (c) names its five vulnerabilities inside the objective itself, so the five
headings below are the objective, not a structure imposed on it.}
\end{cikeybox}

% ------------------------------------------------------------------
\lo{CI-2 a}{Explain key developments in AI since 2017, differentiating between supply-side drivers and demand-side drivers.}

\subsection*{Overview of recent AI developments}

Artificial Intelligence (AI) has been used in financial services for years, but its adoption
has accelerated significantly due to recent technological advancements.

AI applications now span into:

\begin{enumerate}[label=\textbf{\alph*)}]
\item operations (e.g., model risk management, capital optimisation),
\item customer interactions (e.g., chatbots, credit scoring),
\item portfolio management, trading, and
\item regulatory / supervisory technologies (RegTech and SupTech).
\end{enumerate}

\term{Generative AI (GenAI)} has further expanded capabilities in areas such as document
processing, fraud detection, customer support, and market analysis. While these developments
create efficiency and innovation, they also introduce new risks and vulnerabilities across
the financial system.

Since the Financial Stability Board report in 2017, advancements in big data, deep learning,
and computing power, along with the launch of ChatGPT in 2022, have significantly accelerated
AI adoption. Evidence suggests that \term{supply-side factors have been more influential than
demand-side factors} in driving this growth.

\begin{cifigblock}
{\centering
\begin{tikzpicture}[
  font=\sffamily\scriptsize,
  sup/.style={draw=FRMRule, fill=PaleGreen, text width=3.9cm, align=left,
              inner sep=4pt, rounded corners=1pt, minimum height=1.05cm},
  con/.style={draw=FRMRule, fill=PaleRed, text width=3.9cm, align=left,
              inner sep=4pt, rounded corners=1pt, minimum height=1.05cm,
              font=\sffamily\scriptsize\itshape},
  dem/.style={draw=FRMRule, fill=PaleNavy, text width=4.6cm, align=left,
              inner sep=4pt, rounded corners=1pt, minimum height=0.78cm},
  arr/.style={-{Stealth[length=4pt]}, FRMGold, line width=0.7pt}
]
\node[sup] (s1) at (0,0) {\textbf{1. Technological.} Embeddings for unstructured data;
  Transformer architecture; GPUs and cheaper compute};
\node[sup, below=2.5mm of s1] (s2) {\textbf{2. Data-related.} Digital interaction growth;
  unstructured data combined with structured};
\node[sup, below=2.5mm of s2] (s3) {\textbf{3. Business model.} Cloud providers;
  pre-trained models; open-source models};
\node[con, right=9mm of s1] (c1) {Constraint: shortage of skilled human capital, so AI
  talent is expensive and hard to acquire};
\node[con, right=9mm of s2] (c2) {Constraint: not enough real data, so synthetic data,
  which raises data quality and model reliability concerns};
\node[con, right=9mm of s3] (c3) {Constraint: cheaper and faster, but accountability and
  explainability suffer};
\draw[arr] (s1) -- (c1);
\draw[arr] (s2) -- (c2);
\draw[arr] (s3) -- (c3);
\node[above=3mm of s1.north west, anchor=south west,
      font=\sffamily\bfseries\scriptsize\color{FRMGreen}]
  {SUPPLY SIDE \ \ \textnormal{\itshape the more influential driver}};

\coordinate (L) at ($(s3.south west)+(-0.25,-1.05)$);
\coordinate (R) at ($(c3.south east)+(0.25,-1.05)$);
\draw[FRMRule, line width=1.2pt] (L) -- (R);
\fill[FRMGold] ($(L)!0.07!(R)$) coordinate (t1) circle (1.7pt);
\fill[FRMGold] ($(L)!0.58!(R)$) coordinate (t2) circle (1.7pt);
\node[below=3pt of t1, align=center, font=\sffamily\scriptsize\color{FRMNavy}]
  {2017\\FSB report};
\node[below=3pt of t2, align=center, font=\sffamily\scriptsize\color{FRMNavy}]
  {2022\\ChatGPT launch};
\node[above left=2pt and 0pt of R, font=\sffamily\scriptsize\color{FRMGrey}] {today};
\draw[FRMRule, line width=0.7pt] (s3.south) -- (s3.south |- L);

\node[dem] (d1) at ($(R)+(-4.9,-2.15)$) {\textbf{1. Profitability.} Revenue growth, cost
  reduction, better risk forecasting};
\node[dem, below=2.5mm of d1] (d2) {\textbf{2. Competition.} Fear of falling behind peers};
\node[dem, below=2.5mm of d2] (d3) {\textbf{3. Regulation.} Compliance efficiency as rules
  expand};
\node[above=3mm of d1.north west, anchor=south west,
      font=\sffamily\bfseries\scriptsize\color{FRMNavy}]
  {DEMAND SIDE \ \ \textnormal{\itshape motives, not capabilities}};
\draw[FRMRule, line width=0.7pt] (d1.north |- L) -- ([yshift=6mm]d1.north);
\end{tikzpicture}
\par}
\figcap{Figure CI-2.1. Both sides of the split, with the two dated anchors the objective
asks you to reason from. Every supply-side driver comes with its own constraint attached,
drawn to its right; the demand-side drivers in the notes carry none. That asymmetry is the
quickest way to keep the two lists apart: supply-side factors are capabilities that arrived
with a cost, demand-side factors are motives.}
\end{cifigblock}

\subsection*{Supply-side factors}

These factors relate to technological progress, data availability, and evolving business
models.

\begin{enumerate}
\item \term{Technological developments.} Advancements in both hardware and software have
played a key role in AI adoption.
  \begin{enumerate}[label=\alph*)]
  \item Improved processing of unstructured data (text, images, audio) using embeddings.
  \item Emergence of Transformer architecture, enabling large language models and GenAI.
  \item Increased computational power (e.g., GPUs), allowing larger datasets and more complex
  models at lower cost.
  \end{enumerate}
  \emph{Constraint}: shortage of skilled human capital, making AI talent expensive and
  difficult to acquire.
\item \term{Data-related developments.} The rise in digital interactions has led to massive
growth in unstructured data. AI systems can now combine traditional structured data with
sources like images and videos.\\
  \emph{Constraints}: modern AI models require vast amounts of data, which may not always be
  available, hence the use of synthetic data, but this raises concerns about data quality and
  model reliability.
\item \term{Business model developments.} Cloud computing and pre-trained models have
transformed how firms adopt AI.
  \begin{enumerate}[label=\alph*)]
  \item Firms increasingly rely on cloud providers for scalable and cost-efficient AI
  solutions.
  \item Pre-trained models reduce the need to build systems from scratch.
  \item Open-source models are gaining popularity due to flexibility, accessibility, and
  improving performance.
  \end{enumerate}
  \emph{Constraints}: these reduce costs and time, but they also create challenges related to
  accountability and explainability.
\end{enumerate}

\subsection*{Demand-side factors}

\begin{enumerate}
\item \term{Profitability}: a) revenue growth through better personalisation, onboarding, and
lower churn, b) cost reduction via automation and productivity gains, and c) improved risk
management through better forecasting of cash flows, losses, and delinquencies.
\item \term{Competition}: fear of falling behind peers drives faster adoption; this leads to
positives and negatives.
  \begin{itemize}
  \item Positive outcomes: innovation, efficiency, lower costs.
  \item Negative outcomes: rushed implementation without proper safeguards.
  \end{itemize}
\item \term{Regulation}: firms are using AI to improve compliance efficiency as data
protection rules, AI guidance, and international standards expand.
\end{enumerate}

% ------------------------------------------------------------------
\lo{CI-2 b}{Identify and analyze existing and emerging AI use cases in financial services, including industry use cases as well as regulatory and supervisory use cases.}

\subsection*{Industry use cases}

Industry use cases have primarily focused on four areas: customers, operations, trading and
portfolio management, and regulatory compliance.

\begin{enumerate}
\item \term{Customer-focused applications.} On the customer side, it helps banks assess
creditworthiness, improve loan preapproval, personalise marketing, match advisors with
clients, and enhance insurance pricing and risk assessment. GenAI has also made chatbots far
more advanced through LLM-based interactions.
\item \term{Operations-focused applications.} In operations and back-office functions, AI
supports model risk management, capital optimisation, code generation, market impact
analysis, and volatility / liquidity risk management. GenAI has further expanded this by
improving search, content creation, legacy code modernisation, transcription, and access to
advanced tools for firms with limited technical staff.
\item \term{Trading and portfolio management.} Here, GenAI helps with trade execution, market
sentiment analysis from text, and automation of investment analysis and decision-making.
\item \term{Regulatory compliance.} AI has become a key RegTech tool for fraud detection,
AML, CFT, sanctions monitoring, tax evasion detection, trade-based money laundering
detection, and faster financial crime reporting.
\end{enumerate}

\subsection*{Regulatory and supervisory use cases}

\begin{itemize}
\item \term{SupTech (use of technology by regulators)} is rapidly expanding, with central
banks using AI for data collection, analysis, and payment system oversight. AI also enables
real-time economic analysis and use of alternative data, with GenAI/NLP supporting research
and text analysis.
\item In supervision, AI helps review large regulatory documents, draft inspection reports,
and automate workflows. Going forward, GenAI may be used for stress testing and simulating
bank runs, while regulators will need to adopt more advanced AI tools to keep pace with
financial institutions.
\end{itemize}

\begin{cifigblock}
{\centering
\begin{tikzpicture}[
  font=\sffamily\scriptsize,
  ind/.style={draw=FRMRule, fill=PaleNavy, text width=2.75cm, align=center,
              minimum height=1.9cm, inner sep=4pt, rounded corners=1pt},
  reg/.style={draw=FRMRule, fill=PaleGold, text width=5.1cm, align=center,
              minimum height=1.6cm, inner sep=4pt, rounded corners=1pt}
]
\node[ind] (i1) at (0,0) {\textbf{Customers}\\[2pt] creditworthiness, preapproval,
  marketing, advisor matching, insurance pricing, chatbots};
\node[ind, right=4mm of i1] (i2) {\textbf{Operations}\\[2pt] model risk, capital
  optimisation, code generation, market impact, volatility and liquidity risk};
\node[ind, right=4mm of i2] (i3) {\textbf{Trading and portfolio}\\[2pt] execution, sentiment
  from text, automated investment analysis};
\node[ind, right=4mm of i3] (i4) {\textbf{Regulatory compliance}\\[2pt] fraud, AML, CFT,
  sanctions, tax evasion, trade-based ML};
\node[above=3.5mm of i1.north west, anchor=south west,
      font=\sffamily\bfseries\scriptsize\color{FRMNavy}]
  {INDUSTRY USE CASES \ \ \textnormal{\itshape the firm runs the tool}};

\coordinate (DL) at ($(i1.south west)+(0,-0.75)$);
\coordinate (DR) at ($(i4.south east)+(0,-0.75)$);
\draw[FRMRule, line width=0.9pt] (DL) -- (DR);
\node[fill=white, inner sep=3pt, font=\sffamily\scriptsize\itshape\color{FRMGrey},
      anchor=west] at ($(DL)+(0.4,0)$)
  {the dividing line is \emph{who operates it}, not what it does};

\node[reg] (r1) at ($(DL)+(3.0,-2.45)$) {\textbf{SupTech}\\[2pt] data collection, analysis,
  payment system oversight, real-time economic analysis, alternative data};
\node[reg, right=5mm of r1] (r2) {\textbf{Supervision}\\[2pt] reviewing regulatory documents,
  drafting inspection reports, automating workflows; \emph{emerging}: stress testing and
  simulating bank runs};
\node[above=3.5mm of r1.north west, anchor=south west,
      font=\sffamily\bfseries\scriptsize\color{FRMGold}]
  {REGULATORY AND SUPERVISORY USE CASES \ \ \textnormal{\itshape the regulator runs the tool}};
\end{tikzpicture}
\par}
\figcap{Figure CI-2.2. The objective asks for both halves and the halves are told apart by
actor, not by activity. Note that ``regulatory compliance'' sits on the \emph{industry} side:
a firm using AI to detect its own AML breaches is RegTech, an industry use case, while a
central bank using AI to oversee payment systems is SupTech. Same technology, opposite side
of the line.}
\end{cifigblock}

% ------------------------------------------------------------------
\lo{CI-2 c}{Assess the potential implications of AI for financial stability, with a focus on the following vulnerabilities: Third-party dependencies and service provider concentration; Market correlations; Cyber risks; Model risk, data, and governance; Other vulnerabilities, such as fraud, disinformation, and misalignment.}

\subsection*{1. Third-party dependencies and concentration}

Financial institutions increasingly rely on cloud providers and external vendors for AI
capabilities due to high costs and complexity of in-house development. This creates
concentration risk, as only a few firms dominate cloud services, GPUs, and AI model
development.

\begin{cifigblock}
{\centering
\renewcommand{\arraystretch}{1.2}
\begin{tabularx}{\textwidth}{@{}>{\raggedright\arraybackslash}X >{\raggedright\arraybackslash}X@{}}
\toprule
\rowcolor{FRMSoft}
{\sffamily\bfseries\color{FRMRed}Drivers of concentration} &
{\sffamily\bfseries\color{FRMGreen}Factors that may reduce concentration over time}\\
\midrule
a) High costs and complexity of AI systems & a) Growth in open-source models\\
b) Vertical integration across the AI supply chain & b) Increased competition among hardware
providers\\
c) Limited availability of advanced computing chips (GPUs) & c) Development of
domain-specific AI models\\
d) Large investments by incumbent firms & \\
\bottomrule
\end{tabularx}
\par}
\figcap{Table CI-2.1. Concentration is not treated as a one-way ratchet in the notes. Four
forces push it up and three push it down, and open-source models appear on \emph{both}
sides of the ledger in this reading: they reduce concentration here, but in CI-1 c reliance
on the same few open-source models is itself the herding concern.}
\end{cifigblock}

Overall, reliance on third parties exposes firms to operational risks, cyber risks, and
supply chain disruptions, even though it may improve efficiency.

\subsection*{2. Market correlations}

AI-driven decision-making can increase correlations across markets as institutions rely on
similar models, datasets, and signals. This creates \term{herding behaviour}, where multiple
participants react in the same way at the same time.

Key risks include:

\begin{enumerate}[label=\textbf{\alph*)}]
\item Amplified volatility and synchronised market movements.
\item Higher probability of flash crashes and liquidity dry-ups.
\item Correlated errors if models are trained on similar datasets.
\end{enumerate}

For example, if multiple funds use similar sentiment-based AI models, a negative news signal
could trigger simultaneous selling across markets. However, AI can also reduce correlations
if used to create differentiated strategies or niche portfolios.

\begin{cifigblock}
{\centering
\begin{tikzpicture}[
  font=\sffamily\scriptsize,
  st/.style={draw=FRMRule, fill=PaleRed, text width=2.35cm, align=center,
             minimum height=0.95cm, inner sep=3pt, rounded corners=1pt},
  mit/.style={draw=FRMRule, fill=PaleGreen, text width=2.9cm, align=left,
              inner sep=3pt, rounded corners=1pt},
  loop/.style={-{Stealth[length=4pt]}, FRMRed, line width=0.8pt}
]
\node[st] (n1) at (0,1.35)    {Similar models, datasets and signals};
\node[st] (n2) at (3.4,1.35)  {Herding: participants react the same way};
\node[st] (n3) at (3.4,-0.55) {Synchronised moves, flash crashes, liquidity dry-ups};
\node[st] (n4) at (0,-0.55)   {Correlated errors and amplified volatility};

\draw[loop] (n1) -- (n2);
\draw[loop] (n2) -- (n3);
\draw[loop] (n3) -- (n4);
\draw[loop] (n4) -- (n1);

\node[mit] (m1) at (8.6,1.30) {Circuit breakers and trading halts};
\node[mit, below=2.5mm of m1] (m2) {Improved model diversity and transparency};
\node[mit, below=2.5mm of m2] (m3) {Continuous monitoring of systemic risk indicators};
\node[mit, below=2.5mm of m3] (m4) {Differentiated strategies and niche portfolios, which
  \emph{reduce} correlation};
\node[above=3mm of m1.north west, anchor=south west, font=\sffamily\bfseries\scriptsize\color{FRMGreen}] {Where the loop is broken};

\draw[-{Stealth[length=4pt]}, FRMGreen, line width=0.7pt] (m1.west) to[out=180,in=20] (n2.east);
\draw[-{Stealth[length=4pt]}, FRMGreen, line width=0.7pt] (m2.west) to[out=180,in=0] (n3.east);
\draw[-{Stealth[length=4pt]}, FRMGreen, line width=0.7pt] (m4.west) to[out=180,in=-20] (n3.east);
\end{tikzpicture}
\par}
\figcap{Figure CI-2.3. The correlation vulnerability is a closed loop, not a chain: correlated
errors feed back into the model similarity that caused them. The three mitigations in the
notes do not all attack the same point, which is why they are not substitutes. Circuit
breakers interrupt the reaction, model diversity interrupts the transmission, and
differentiated strategies are the one item in the reading that runs the loop backwards.}
\end{cifigblock}

\subsection*{3. Cyber risks}

AI significantly increases both the scale and sophistication of cyber threats. GenAI enables
attackers to automate and enhance attacks while lowering entry barriers for less-skilled
actors.

Key vulnerabilities include:

\begin{enumerate}[label=\textbf{\alph*)}]
\item Prompt injection attacks (extracting confidential data).
\item Data and model poisoning (manipulating training data).
\item Increased attack vectors due to large datasets and system complexity.
\end{enumerate}

Critical financial entities (e.g., central banks, payment systems, major cloud providers) are
especially vulnerable. A successful attack can lead to operational failures, loss of
confidence, and market instability.

\begin{citrapbox}[title={An asymmetry, stated in the notes and worth holding exactly as written}]
While AI strengthens cybersecurity defences (e.g., faster detection, automated response),
malicious actors may benefit more in the short term, creating an asymmetry in risk. The claim
is not that AI helps attackers and not defenders. It is that it helps both, and helps
attackers \emph{more}, and \emph{in the short term}. Both qualifiers are load-bearing.
\end{citrapbox}

\subsection*{4. Model risk, data, and governance}

As AI models become more complex, model risk increases, especially during stressed market
conditions.

Challenges include:

\begin{enumerate}[label=\textbf{\alph*)}]
\item Low explainability (``black box'' models).
\item Difficulty validating outputs, especially unstructured text from LLMs.
\item Data quality issues due to reliance on large, unstructured datasets.
\item Model hallucinations (confident but incorrect outputs).
\end{enumerate}

To counter this, strong governance is essential, including data quality checks, model
validation and performance monitoring, and continuous oversight and risk mitigation
strategies.

\subsection*{5. Other vulnerabilities}

\begin{enumerate}
\item \term{Financial fraud.} GenAI enables advanced fraud techniques such as deepfakes, fake
identities, and bypassing security systems. Fraud may become easier to execute than detect,
especially as malicious actors adopt AI faster than regulated institutions. Improving AI
literacy and financial awareness is key to mitigating these risks.
\item \term{Disinformation.} AI can amplify misinformation through scalable and convincing
content generation, especially via social media. This can trigger strong market reactions and
increase financial instability during stressed conditions.
\item \term{Misalignment.} AI misalignment occurs when systems optimise profits but act
against regulatory or ethical standards, leading to risks like collusion or disinformation
that can destabilise markets.
\end{enumerate}

\begin{cikeybox}[title={Closing position of the reading}]
In the long run, AI increases demand for data, computing, and energy, while creating uneven
regulation across firms. It may also impact inflation, interest rates, and jobs, with even
high-skilled workers at risk. Although AI improves efficiency, compliance, and analytics, it
also introduces new risks due to rapid adoption and limited data.\par\vspace{4pt}
Regulators should focus on: closing data gaps; updating regulatory frameworks; strengthening
supervision.\par\vspace{4pt}
Overall, AI brings strong benefits but also growing systemic risks that require careful
oversight.
\end{cikeybox}
